% Source: https://artofproblemsolving.com/community/c3h2861971_simple_geoalg_combination
% Stable identity: aops-topic-2861971-q01
The right angles of the rectangle subtend its diagonals, so each diagonal is a diameter, of length $2$.
\begin{asy}
usepackage("amsmath");
size(3.5cm);
defaultpen(fontsize(10pt));
real a=degrees(atan2(1,1.6));pair O=(0,0);pair P=dir(a),Q=dir(180-a),R=dir(180+a),S2=dir(-a);
draw(circle(O,1));draw(P--Q--R--S2--cycle);draw(P--R,dashed);dot(O);label("$O$",O,S);
label("$a$",(P+Q)/2,N);label("$b$",(P+S2)/2,E);label("$2$",(0.35,0.3),NW);
\end{asy}

With sides $a,b$: $a^2+b^2=4$. By the QM-AM inequality (or $(a-b)^2\ge0$), $(a+b)^2\le2(a^2+b^2)=8$, so $a+b\le2\sqrt2$, with equality when $a=b=\sqrt2$ (a square). The maximum perimeter is $2(a+b)=\boxed{4\sqrt2}$.
